\subsection{Whole-domain retained square moments}
\label{sec:unit-retained-whole-moments}
\begin{lemma}[Whole-domain data for $\sigma=2$]
\label{lem:unit-retained:whole-data}
Take the actual $U$ of Lemma~\ref{lem:unit-retained-pair} with
$\sigma=2$, and put $Z=KU$. The following data hold for the
complete interval $0\le d\le1/3$:
\[
 \|U\|^2\le d\,v_n(d)^2,\quad
 \E U\le d\,z_n(d),\quad \E Z\ge\alpha_n(d),\quad
 |Z|\le v_n\sqrt d,\quad \E Z^2\le\theta d\,v_n^2,
\]
where
\[
 \begin{array}{c|ccc}
 n&v_n&z_n&\alpha_n\\ \hline
 5&3/8+23d/16&\theta(1+d)&(51/32+d)d^2\\
 \ge6&3/8+3d/2&79/256+13d/16+d^2&(51+19d)d^2/32.
 \end{array}
\]
\end{lemma}
\begin{proof}
The original spectral measure at $u$ has mass $d$ and
$\int x=\delta_2$. Skew adjunction gives
$\|U\|^2=\int(a^2+xb^2)$, $\E U=\int b$,
and $\E KU=\int a$.
For $n=5$, $a'=1-2x>0$ and $b'=1$. First maximize the coordinates in $x$ and
then use $\delta_2\le\theta d$. The constant-vector squared norm
is $131313/1048576<9/64$ and the linear-vector squared norm is
$66997/32768<(23/16)^2$; the latter margin is $715/32768$.
Also $b\le\theta(1+d)$ and
\[
 \E Z=d^2+(1+d)\delta_2-\delta_3
       \ge(51/32+d)d^2.
\]

For $n\ge6$ put $u_\theta=10101/32768$.
The exact coordinates are bounded above by
\[
 x-x^2+x^3+d+\delta_2+E_0(d-x),\qquad
 x-x^2+x^3+2\delta_2+dE_0.
\]
They increase in $x$, with derivatives at least $1/3$ and $2/3$.
At $x=\theta$, then $\delta_2\le\theta d$, the bounding vector is
\[
 (u_\theta+(1+\theta)d+E_0(d-\theta),
   \sqrt\theta[u_\theta+2\theta d+dE_0]).
\]
Its squared norm is convex in $E_0\in[0,d]$.
At $E_0=0$, the constant norm is below $3/8$ and the linear
squared norm is $18397/8192<9/4$. At $E_0=d$, increase
$u_\theta$ to $u_*=79/256$.
The difference
\[
 (3/8+3d/2)^2-(u_*+d+d^2)^2
                  -\theta(u_*+2\theta d+d^2)^2
\]
has degree four and Bernstein coefficients on $[0,1/3]$
\[
 \frac{14067}{2097152},\quad\frac{67211}{2097152},\quad
 \frac{1123003}{18874368},\quad\frac{1216019}{18874368},\quad
 \frac{30457}{6291456}.
\]
All are positive. This proves the norm bound at both $E_0$
endpoints and hence for the complete interval.
Further $b\le u_\theta+2\theta d+d^2<z_n$.
Using $\beta_n\le d-\kappa\delta_2$, $a_N\ge1-x$,
$E_0\le d$, and $\delta_3\le\theta\delta_2$ gives
\[
 \begin{split}
 \int a&\ge d^2+\kappa(1+d)\delta_2-E_0(\delta_2-d^2)\\
        &\ge d^2+(\kappa-\theta d)\delta_2+d^3
         \ge(1+\kappa+\kappa d)d^2.
 \end{split}
\]
The last coefficient is positive throughout the full domain.
Finally apply the original row and spectral bounds to $KU$.
At $d=0$ every vector appearing here vanishes.
\end{proof}

\begin{proposition}[One host-domain payment for every $(2,n)$, $n\ge5$]
\label{prop:unit-retained:whole-L1}
On $0\le d\le1/3$, in the same canonical direction,
\[
 H=pT(h_2f_n),\quad U_{\rm out}=T^2H
 \qquad\Longrightarrow\qquad
 I_q(U_{\rm out})^3\le e^{-d/500}.
\]
\end{proposition}
\begin{proof}
Put $m=1-d$, $v=v_n$, $z=z_n$, $\alpha=\alpha_n$.
The actual square in \eqref{eq:unit-retained:dual-sums} gives
\[
 T(h_2f_n)\le m^2q+m(\E U-KU)+\|U\|^2/2.
\]
AM--GM between the two nonnegative original factors gives
\[
 0\le H\le(c-Z/2)^2,\qquad
 c=1-d+dz/2+dv^2/[4(1-d)].
\]
Here $c\ge2/3>7/16\ge v\sqrt d/2$.
Lemma~\ref{lem:unit-square-hermite}
therefore applies to the actual $X=-Z$, with cap $v\sqrt d$,
mean upper bound $-\alpha$, and raw second bound
$\theta d v^2$. Its moment compatibility follows from these same
actual moments.

Let $d=t^2$, $a=\alpha/d$, $u=v$, and use the polynomial
definitions $P,Q,L,B,N,\mathcal D$ in
Section~\ref{sec:unit-square-payments}. Take
\[
 \xi=\tfrac12(24427/3420+1/500)d-\tfrac52d^2,
 \qquad\mathcal T=\sum_{j=0}^7(-\xi)^j/j!.
\]
The exact factor in both rows is $t^2$:
$\mathcal G=(\mathcal T\mathcal D-N)/t^2$.
For $n=5$ take
\[
 C=1024-1804d+1300d^2+321d^3,\qquad D=1024(1-d);
\]
for $n\ge6$ take
\[
 C=512-927d+785d^2+336d^3-256d^4,\qquad D=512(1-d).
\]
These are the actual $c=C/D$ above. On the single interval
$0\le t\le577351/10^6$, the two polynomials have degrees
$85,99$, respectively. Their left values exceed $1/2$;
after division by their left values, every one of their
$86+100=186$ Bernstein coefficients is at least $9/10$.

Lemma~\ref{lem:bernstein-certificate} proves
$0\le\E H^4\le\mathcal T\le e^{-\xi}$; the last inequality is
the odd Taylor lower bound, since $\xi\ge0$.
With the ordinary-fourth output factor of paragraph~\ref{part:R29.11},
$B_2\le e^{(24427/3420)d-5d^2}$, this gives
$I_q(U_{\rm out})^3\le B_2(\E H^4)^2\le e^{-d/500}$.
At $d=0$ equality is evaluated directly.
\end{proof}
